\chapter{Physikalisches E-Bike-Modell}
\label{ch:physik}

Das Fahrzeugmodell in \texttt{ebike.py} berechnet fuer jeden GPS-Punkt die benoetigte Kraft aus Fahrwiderstaenden, Beschleunigung und Steigung. Die Rechnung erfolgt in Fahrtrichtung und nutzt die aus dem GPS-Datensatz abgeleitete Geschwindigkeit und Beschleunigung.

\section{Einzelkraefte}

Der Luftwiderstand folgt im Projekt direkt der bekannten quadratischen Form und entspricht der in der Aerodynamik ueblichen Drag-Gleichung \cite{nasa_drag}:
\begin{equation}
F_{\mathrm{Luft}} = \frac{1}{2}\rho c_{\mathrm{w}}A v^2
\end{equation}
Dabei ist \(\rho\) die Luftdichte, \(c_{\mathrm{w}}A\) die aerodynamische Kennzahl und \(v\) die Geschwindigkeit. Fuer die Luftdichte verwendet das Projekt entweder die ueber \texttt{luftdichte.py} berechnete aktuelle Dichte oder eine konstante Standarddichte, falls Hoehe und Temperatur nicht nutzbar sind.

Die Hangabtriebskraft lautet
\begin{equation}
F_{\mathrm{Hang}} = m g \sin(\varphi)
\end{equation}
und die Beschleunigungskraft
\begin{equation}
F_{\mathrm{Beschleunigung}} = m a
\end{equation}
mit der Gesamtmasse \(m\), der Erdbeschleunigung \(g\), der Steigung \(\varphi\) und der Beschleunigung \(a\).

Der Rollwiderstand wird im Code als
\begin{equation}
F_{\mathrm{Roll}} = c_{\mathrm{R}} m g \cos(\varphi)
\end{equation}
modelliert. Die verwendete Masse setzt sich aus Fahrer- und Fahrradmasse zusammen.

\section{Gesamtkraft und Leistung}

Die Funktion \texttt{EBike.antriebskraft\_N()} bildet die Summe
\begin{equation}
F_{\mathrm{Antrieb}} = F_{\mathrm{Luft}} + F_{\mathrm{Hang}} + F_{\mathrm{Beschleunigung}} + F_{\mathrm{Roll}}
\end{equation}
und damit genau die im Quellcode verwendete Kraft. Aus Kraft und Geschwindigkeit folgt die mechanische Leistung
\begin{equation}
P = F v
\end{equation}
und das Raddrehmoment
\begin{equation}
M = F r
\end{equation}
mit dem Radradius \(r\).

\section{Parameter}

Die Parameter werden aus \texttt{data/bike\_config.yaml} geladen. Aktuell verwendet das Projekt eine Fahrer- und Radmasse von zusammen \SI{80}{\kilogram}, eine aerodynamische Kennzahl von \SI{0.5625}{\square\metre}, einen Raddurchmesser von \SI{27}{\centi\metre} und einen Rollwiderstandskoeffizienten von \num{0.005}. Diese Werte bestimmen gemeinsam mit der GPS-Kinematik die Leistungsanforderung entlang der Strecke.
